\section{Preliminaries: What Must Be Fixed}
\label{sec:preliminaries}

This section is presently a roadmap.  The final paper should reproduce every
definition whose details affect the YAMSO claim, while importing routine
cryptographic syntax and standard security definitions by reference.

\subsection{Conventions}

We write $\lambda$ for the security parameter, $[n]=\{1,\ldots,n\}$, and
$\bits^*$ for the set of finite bit strings.  All adversaries and protocols
are probabilistic polynomial time unless explicitly stated otherwise, and
$\negl(\lambda)$ denotes a negligible function.  Protocol identifiers,
epochs, wire identifiers, public-key indices, ledger positions, and object
types must be included in every cryptographic domain separator.  A computation
key is assigned to exactly one canonical protocol role.

\subsection{Definitions that belong in this paper}

The following objects are model-specific and should eventually be defined in
full.

\begin{description}
  \item[Ledger.] One synchronous, append-only list $L$ seen identically by
    every party.  It begins with a generic CRS, grows by registration and
    scheduled message entries, and performs no protocol-specific validation.
    It has no cryptographic certification key.

  \item[Growing key registry.] An adversarial registration request creates a
    computation PID equal to its fresh public key.  A registration from an
    existing application PID creates and publicly binds a protected key to
    that named party.  Neither request contains a protocol role; registration
    order determines roles at the protocol level.

  \item[YAMSO turn.] A public initialization parameter
    $\mathsf{WhosTurn}$ selects one unspent PID from the current $L$.  The
    selected party either produces one message or receives an explicit
    $\mathsf{NoMsg}$ entry and is then spent forever.

  \item[Protected local evaluation.] An honest scheduled party supplies an
    algorithm $A$ privately to the ledger, which evaluates $A(sk,L)$ without
    first returning $sk$ to the party.  This lets an ordinary role reveal its
    entire key and an input role reveal only an input-dependent function of
    its state through the same generic interface.

  \item[Two corruption rules.] Each computation PID is independently corrupt
    with probability $p$ when registered and, if honest, cannot later be
    corrupted.  Named application parties have the ordinary adaptive
    corruption interface in the model.  The first construction statement
    restricts their corrupt set to be fixed before setup.

  \item[Closed input slot.] At a protocol-defined cutoff, the public ledger
    prefix determines one immutable input record.  Public branch-indexed
    one-way tags determine whether it contains exactly one valid activation
    token.  A missing, ambiguous, or malformed record selects the structural
    default branch; late messages are ignored.

  \item[Common-input certificate.] Fresh one-use roles hold opposite one-way
    preimages for the coordinates of an encoded digest.  Honest roles compute
    the value to be certified from the same closed ledger cut, rather than
    signing a caller-supplied string.  Verification permits enough missing
    openings for liveness but must rule out a second accepted digest.

  \item[OGCS prefix.] A fixed finite public object is evaluated on complete
    finite histories.  Outputs on every accepting history are literal
    prefixes.  A history that first differs at some accepted command must not
    reveal the hidden logical state of the canonical history; a rejected
    history never extends again.
\end{description}

\subsection{Primitives that can largely be imported}

Standard definitions of pseudorandom functions, collision-resistant and
one-way hashing, commitments, secret sharing, public-key encryption, and
circuit garbling need not be repeated in full.  We should state the exact
variants used and cite them.  The following interfaces require particular
care.

\begin{description}
  \item[Role-state encryption.] Role membership is public in the preliminary
    model, so ordinary multi-user encryption can protect the state placed
    under each registered key; recipient anonymity is not required by this
    abstraction.  The semantic association between the two structural
    committees for an input wire must nevertheless remain hidden.  Each key
    has one canonical role packet.  Off-path evaluations may address further
    packets to the same key, and later disclosure decrypts them; the OGCS
    isolation lemma must make all such plaintexts independently rooted fork
    material.

  \item[Verifiable threshold sharing.] The construction needs privacy below
    the reconstruction threshold, reconstruction from honest shares, and
    public rejection of malformed shares or keys.  Public verification data
    must not disclose the semantic branch permutation.

  \item[Two-stage garbled circuits.] Preparation fixes external-input label
    pairs, tokens, and role packets before inputs arrive.  Completion occurs
    only after input closure and must support prescribed active external and
    carried-state labels, prescribed public output, hidden next-state label
    pairs, and consistency with every later-exposed selected role.  This
    causal interface is stronger than a bare one-shot privacy statement.

  \item[High-distance one-time opening.] The certificate uses the distance
    idea from Lamport-style leakage-resilient signatures
    \cite{KatzVaikuntanathan2009LeakageSignatures}.  The paper must prove its
    own partial-opening liveness and non-equivocation statement under random
    role corruption and complete post-speech exposure, including an explicit
    code-rate-versus-tail calculation.

  \item[Succinct TM-iO and puncturing.] For a declared input-length bound, the
    candidate OGCS can use a compact obfuscated Turing-machine program with
    puncturable PRF/KDF state
    \cite{KoppulaLewkoWaters2014IOTM,AnanthJainSahai2017TMIO}.  We must state
    the bound, succinctness, functionality-equivalence condition, and any
    equal-running-time requirement.  The cited general-purpose
    ordinary-iO-based bounded-input results do not by themselves realize the
    horizon-free replay interface of \cref{sec:ongoing}.

  \item[Delayed backdoor programming.] Hofheinz et al.'s adaptive universal
    sampler provides the hidden-trigger paradigm
    \cite{HofheinzEtAl2014UniversalSamplers}.  Selectively secure pseudorandom
    puncturable deterministic encryption gives a related formulation
    \cite{KoppulaPoelstraWaters2017FastSamplers}.  Their theorems do not apply
    automatically to our stateful syntax; the OGCS replacement lemma must be
    proved separately.

  \item[Statistical binding for iO hybrids.] Computational unforgeability does
    not by itself justify iO between programs that differ on an invalid but
    mathematically possible input.  In its binding setup, a somewhere
    statistically unforgeable signature makes alternate valid signatures at a
    selected round and message statistically nonexistent, except with
    negligible setup probability \cite{FernandoEtAl2023SSU}.  Adapting that
    property to our certificate chain is an open proof obligation.
\end{description}

\paragraph{Proof boundary.}
The main text uses the generic ledger interface of \cref{sec:model}.  Once the
API is stable, a detailed $\mathcal F_{\mathsf{Ledger}}$ can be placed in an
appendix.  The target functionality contains only named application parties;
the randomly corrupted computation population is implementation machinery to
be simulated away.  Intermediate cryptographic components are handled by
whole-process replacement lemmas, avoiding hybrid interfaces that export
their secret trapdoors to the environment.
